\section{The Brownian--Poisson MV-FBSDEJ}\label{sec:main-assumptions}
Let $(\Omega,\F,(\F_t),\Prob)$ be a filtered probability space satisfying the usual
conditions, over a fixed horizon $T>0$ and dimensions $n,m,d\ge1$, carrying a
$d$-dimensional Brownian motion $W$ and an independent Poisson random measure $N$ on
$\R_+\times E$ with intensity $\diff t\,\nu(\diff e)$, where $\nu$ is a
$\sigma$-finite measure on a countably generated $(E,\mathcal E)$ (so that
$L^2(\nu;\R^j)$ is separable for every $j\ge1$). The filtration $(\F_t)$ is the usual augmentation of the one generated by $W$,
$N$ and an independent initial $\sigma$-field $\F_0$.
We do not require $\nu$ to be finite: the infinite-activity regime,
in which $N$ has infinitely many atoms on every nondegenerate time interval, is admitted
throughout, with the finite-activity case $\nu(E)<\infty$ recovered as a special instance.

Write $\Nt(\diff t,\diff e):=N(\diff t,\diff e)-\diff t\,\nu(\diff e)$ for the
associated compensated measure. For predictable $\psi\ge0$ the compensation formula
reads
\begin{equation}\label{eq:compid}
\E\!\int_0^T\!\!\int_E\psi_s(e)\,N(\diff s,\diff e)
=\E\!\int_0^T\!\!\int_E\psi_s(e)\,\nu(\diff e)\,\diff s\in[0,\infty],
\end{equation}
and the zero-mean property of the compensated integral reads
\begin{equation}\label{eq:zeromean}
\E\!\int_0^T\!\!\int_E\psi_s(e)\,\Nt(\diff s,\diff e)=0
\qquad\text{whenever}\quad
\E\!\int_0^T\!\!\int_E|\psi_s(e)|\,\nu(\diff e)\,\diff s<\infty .
\end{equation}

The unknown is the tuple $\Theta=(X,Y,Z,U)$, with law $\mu_t=\Law(\Theta_t)$ and predictable
counterparts $\Theta_t^-,\mu_t^-$ introduced below; the continuous MV-FBSDEJ is the system
\eqref{eq:mvfbsdej} displayed in the introduction. The tuple
$\Theta_t=(X_t,Y_t,Z_t,U_t)$ takes values in the Hilbert space $\mathcal H$,
componentwise
\[
X_t\in\R^n,\qquad Y_t\in\R^m,\qquad Z_t\in\R^{m\times d},\qquad U_t\in L^2(\nu;\R^m).
\]
The jump terms are evaluated at the predictable tuple
\[
\Theta_t^-=(X_{t-},Y_{t-},Z_t,U_t),\qquad \mu_t^-=\Law(\Theta_t^-),
\]
in which the forward and backward states are replaced by their left limits while the
integrands $Z_t,U_t$, already predictable, are left untouched.
Throughout, $\mu_t$ and $\mu_t^-$ denote the Borel
representatives fixed in Lemma~\ref{lem:lawflow}, which agree with $\Law(\Theta_t)$ and
$\Law(\Theta_t^-)$ for a.e.\ $t$. Expressions such as $b(t,\Theta_t,\mu_t)$ are thereby
defined at every $t$, while every claim made about them is $\diff t$-a.e.\ or
time-integrated.

The filtration generated by $W$ and $N$ has the \emph{Brownian--Poisson predictable
representation property}~(\cite[\S5.3]{Applebaum}; \cite[Theorem~13.49]{HeWangYan}), and it
passes to the independent initial enlargement by $\F_0$ in the standard way:
every square-integrable $(\F_t)$-martingale $M$ can be written as
\[
M_t=M_0+\int_0^t Z_s\,\diff W_s+\int_0^t\!\!\int_E U_s(e)\,\Nt(\diff s,\diff e),
\qquad\text{for every }t\in[0,T],
\]
for an $\F_0$-measurable $M_0$ and a predictable square-integrable pair $(Z,U)$, and this
pair is unique up to null sets.

\begin{definition}[Diagonal input]\label{def:diag}
Fix $t\in[0,T]$. A \emph{diagonal input at time $t$} is a pair $(\Theta,\mu)$ consisting of a
square-integrable $\F_t$-measurable tuple $\Theta\in L^2(\Omega,\F_t;\mathcal H)$, on the
standing basis, \emph{together with its own law} $\mu=\Law(\Theta)\in\mathcal P_2(\mathcal H)$;
such pairs form the diagonal
\[
\mathcal D_t:=\bigl\{\bigl(\Theta,\Law(\Theta)\bigr):\Theta\in L^2(\Omega,\F_t;\mathcal H)\bigr\}
\ \subset\ L^2(\Omega,\F_t;\mathcal H)\times\mathcal P_2(\mathcal H).
\]
\end{definition}

\noindent All hypotheses below are imposed along $\mathcal D_t$ only: the measure argument is
never varied independently of the tuple. Throughout, a lower-case $\theta$ denotes a \emph{point} of
$\mathcal H$ and an upper-case $\Theta$ denotes a \emph{random} $\mathcal H$-valued tuple.

\smallskip\noindent\textbf{Difference notation.}
For two diagonal inputs $(\Theta,\mu),(\Theta',\mu')$ at the same time $t$ on the same
probability space, $\delta$ denotes the difference
\[
\delta X:=X-X',\qquad
\delta\psi:=\psi(t,\Theta,\mu)-\psi(t,\Theta',\mu'),\quad\psi\in\{b,\sigma,h,f\},
\]
and, with the \emph{induced} terminal laws $\rho:=\Law(X_T)$ and
$\rho':=\Law(X'_T)$,
\[
\delta g_T:=g(X_T,\rho)-g(X'_T,\rho') .
\]

Throughout, $b,\sigma,h,f,g$ are deterministic functions of their arguments; randomness
enters only through the inputs and, in Section~\ref{sec:wp-proof}, through additive
square-integrable forcings (the $\omega$-random coefficients of~\cite{ChenNie,LiMin} are
not pursued here). Their signatures are
\[
\begin{aligned}
b&\colon\ [0,T]\times\mathcal H\times\mathcal P_2(\mathcal H)\ \to\ \R^n,\\
\sigma&\colon\ [0,T]\times\mathcal H\times\mathcal P_2(\mathcal H)\ \to\ \R^{n\times d},\\
h&\colon\ [0,T]\times\mathcal H\times\mathcal P_2(\mathcal H)\ \to\ L^2(\nu;\R^n),\\
f&\colon\ [0,T]\times\mathcal H\times\mathcal P_2(\mathcal H)\ \to\ \R^m,\\
g&\colon\ \R^n\times\mathcal P_2(\R^n)\ \to\ \R^m,
\end{aligned}
\]
the terminal function alone carrying no time argument and depending on the state and its
law. All coefficients are jointly Borel measurable in their arguments; for the jump
coefficient, whose values lie in the fibre, this is taken as Borel measurability of the map
\begin{equation}\label{eq:hborel}
(t,\theta,\mu)\ \longmapsto\ h(t,\theta,\cdot,\mu),
\qquad
[0,T]\times\mathcal H\times\mathcal P_2(\mathcal H)\ \longrightarrow\ L^2(\nu;\R^n).
\end{equation}
Since $L^2(\nu;\R^n)$ is separable, \eqref{eq:hborel} follows from joint pointwise
measurability of $(t,\theta,e,\mu)\mapsto h(t,\theta,e,\mu)$: the latter makes $h$ weakly
measurable by Tonelli, and the Pettis measurability
theorem~\cite[Theorem~1.1.6]{HNVW} then upgrades weak to strong measurability.

For predictable processes we use the Fubini isometry
\begin{gather}
\mathcal P:=\text{predictable $\sigma$-field on }\Omega\times[0,T]\text{ under }\diff t\otimes\diff\Prob,\notag\\
L^2(\mathcal P;\Lnu)\ \cong\ L^2(\mathcal P\otimes\mathcal E)\label{eq:fubiniiso}
\end{gather}
of~\cite[Propositions~1.2.24 and~1.2.25]{HNVW}, which supplies a
$\mathcal P\otimes\mathcal E$-measurable representative, unique up to
$\diff t\otimes\diff\Prob\otimes\nu$-null sets; every Poisson integral below, and every
square-integrable $\Lnu$-valued predictable integrand appearing in one, is read through
such a representative. In particular, a pointwise formula $u(e)$ for $u\in\Lnu$ is
meaningful only when it descends to a well-defined element of $\Lnu$, independent of
the representative of $u$. Consequently, along any diagonal input,
\begin{equation}\label{eq:coefmeas}
\begin{aligned}
(s,\omega)&\mapsto\psi\big(s,\Theta_s(\omega),\mu_s\big),\ \ \psi\in\{b,\sigma,f\},
&&\text{progressively measurable},\\[2pt]
(s,\omega,e)&\mapsto h\big(s,\Theta_{s-}(\omega),e,\mu_s^-\big),
&&\text{predictable}.
\end{aligned}
\end{equation}
The well-definedness of the law flows $t\mapsto\mu_t,\mu_t^-$ and the independence of all
time-integrated coefficient identities from the choice of predictable representatives of
$Z,U$ are established in Lemma~\ref{lem:lawflow} in Appendix~\ref{app:lawflow}.

\begin{assumption}[Integrability at the origin]\label{ass:integ}
$\xi_0\in L^2(\Omega,\F_0;\R^n)$ and, with $\delta_0$ the Dirac mass at the relevant origin,
the origin values satisfy
\[
\begin{aligned}
b(\cdot,0,\delta_0)\ &\in\ L^2(0,T;\R^n),
&\qquad\quad \sigma(\cdot,0,\delta_0)\ &\in\ L^2(0,T;\R^{n\times d}),\\[2.2ex]
h(\cdot,0,\cdot,\delta_0)\ &\in\ L^2\bigl(0,T;L^2(\nu;\R^n)\bigr),
&\qquad\quad f(\cdot,0,\delta_0)\ &\in\ L^2(0,T;\R^m).
\end{aligned}
\]
\end{assumption}
\begin{assumption}[Lipschitz coefficients]\label{ass:lip}
There is $L\ge0$ such that, for a.e.\ $t\in[0,T]$ and all
$(\Theta,\mu),(\Theta',\mu')\in\mathcal D_t$ (Definition~\ref{def:diag}),
\[
|\delta b|+|\delta\sigma|_F+\Lnorm{\delta h}+|\delta f|
\ \le\ L\bigl(|\delta X|+|\delta Y|+|\delta Z|_F+\Lnorm{\delta U}+W_2(\mu,\mu')\bigr)
\]
almost surely; and, for all $X_T,X'_T\in L^2(\Omega,\F_T;\R^n)$,
\begin{equation}\label{eq:lipT}
|\delta g_T|\ \le\ L\bigl(|X_T-X'_T|+W_2(\rho,\rho')\bigr)
\end{equation}
almost surely.
\end{assumption}

\begin{assumption}[Coupling matrix]\label{ass:G}
$G\in\R^{m\times n}$ has full rank. Write $\|G\|$ for its operator norm, so that
$|Gv|\le\|G\|\,|v|$ and $|\Gstar v|\le\|G\|\,|v|$, and $s_{\min}(G)>0$ for its
smallest singular value; set
\[
c_G:=s_{\min}(G)^{-2} .
\]
Then
\[
|x|^2\le c_G|Gx|^2\quad(m\ge n),
\qquad\qquad
|y|^2\le c_G|\Gstar y|^2\quad(n\ge m).
\]
When $\Gstar$ is injective the second bound extends to the matrix- and $\Lnu$-valued
arguments,
\[
|Z|_F^2\le c_G|\Gstar Z|_F^2,
\qquad\qquad
\Lnorm{U}^2\le c_G\Lnorm{\Gstar U}^2 .
\]
\end{assumption}
\begin{assumption}[Terminal and interior jump-extended $G$-monotonicity]\label{ass:mono}
For all $X_T,X'_T\in L^2(\Omega,\F_T;\R^n)$,
\[
\E\inner{\delta g_T}{G\,\delta X_T}\ge\alpha\,\E|G\delta X_T|^2,
\]
and, for a.e.\ $t\in[0,T]$ and all $(\Theta,\mu),(\Theta',\mu')\in\mathcal D_t$,
\begin{multline}\label{eq:mono}
\E\!\Big[\inner{\delta b}{\Gstar\delta Y}+\inner{\delta\sigma}{\Gstar\delta Z}_F
+\inner{-\delta f}{G\delta X}+\!\int_E\!\inner{\delta h(e)}{\Gstar\delta U(e)}\nu(\diff e)\Big]\\
\le-\beta_1\E\big[|\Gstar\delta Y|^2+|\Gstar\delta Z|_F^2+\Lnorm{\Gstar\delta U}^2\big]-\beta_2\E|G\delta X|^2
\end{multline}
The constants satisfy
\[
\alpha,\beta_1,\beta_2\ge0,\qquad \alpha+\beta_1>0,\qquad \beta_1+\beta_2>0,
\]
with $\beta_1>0$ if $m<n$ and $\alpha,\beta_2>0$ if $m>n$.
\end{assumption}
\begin{remark}[Wasserstein bounds]\label{rem:wasserstein-bounds}
The pair $(\Theta,\Theta')$ realises a coupling of $(\mu,\mu')$ on the common probability
space, so the $W_2$-infimum~\cite[Definition~6.1]{Villani2009} is bounded by its synchronous transport
cost:
\begin{align}
W_2(\mu,\mu')^2&\le\E\big[|\delta X|^2+|\delta Y|^2+|\delta Z|_F^2+\Lnorm{\delta U}^2\big],
\label{eq:coupling}\\
W_2(\rho,\rho')^2&\le\E|\delta X|^2 ,\label{eq:couplingT}
\end{align}
\eqref{eq:couplingT} being the $X$-marginal case of \eqref{eq:coupling}.

In the other direction, let $K$ be a separable Hilbert space and $\bar\phi:\mathcal H\to K$
Lipschitz. The law functional $\mu\mapsto\int_{\mathcal H}\bar\phi\,\diff\mu$ is then
$W_2$-Lipschitz with the same constant: writing the difference as an integral against an
optimal coupling $\pi$ of $(\mu,\mu')$, the Lipschitz bound and Jensen give
\begin{align}
\Big|\int_{\mathcal H}\bar\phi\,\diff\mu-\int_{\mathcal H}\bar\phi\,\diff\mu'\Big|_K
&=\Big|\int_{\mathcal H\times\mathcal H}\big(\bar\phi(\theta)-\bar\phi(\theta')\big)\,\diff\pi\Big|_K\nonumber\\
&\le\mathrm{Lip}(\bar\phi)\int_{\mathcal H\times\mathcal H}|\theta-\theta'|_{\mathcal H}\,\diff\pi\nonumber\\
&\le\mathrm{Lip}(\bar\phi)\Big(\int_{\mathcal H\times\mathcal H}|\theta-\theta'|_{\mathcal H}^2\,\diff\pi\Big)^{1/2}\nonumber\\
&\le\mathrm{Lip}(\bar\phi)\,W_2(\mu,\mu') .\label{eq:lipfun}
\end{align}
Marginal projection $\mathcal P_2(\mathcal H)\to\mathcal P_2(\Lnu)$ is $1$-Lipschitz for $W_2$
by the same argument: write $\mu^U\in\mathcal P_2(\Lnu)$ for the $U$-marginal of $\mu$. The image
of this optimal $\pi$ under $(\theta,\theta')\mapsto(u,u')$ is a coupling of
$(\mu^U,\mu'^U)$, so bounding $W_2(\mu^U,\mu'^U)$ by its cost gives
\begin{align}
W_2(\mu^U,\mu'^U)^2
&\le\int_{\mathcal H\times\mathcal H}\Lnorm{u-u'}^2\,\diff\pi\nonumber\\
&\le\int_{\mathcal H\times\mathcal H}|\theta-\theta'|_{\mathcal H}^2\,\diff\pi\nonumber\\
&\le W_2(\mu,\mu')^2 .\label{eq:margproj}
\end{align}
\end{remark}
\begin{remark}[The coupling matrix $G$]\label{rem:G-pairing}
The full-rank matrix $G\in\R^{m\times n}$ serves to pair the forward and backward
components when $n\neq m$; its full-rank property ensures the resulting $G$-coercivity
controls the full forward variables when $m\ge n$, or the full backward variables when
$n\ge m$.
\end{remark}
\begin{remark}[Expected diagonal monotonicity]\label{rem:diag-mono}
Assumption~\ref{ass:mono} is imposed \emph{in expectation}. The \emph{pointwise}
$G$-monotonicity condition requires the inequality of \eqref{eq:mono}, without expectations,
at every fixed $\theta,\theta'\in\mathcal H$ and $\mu,\mu'\in\mathcal P_2(\mathcal H)$.
Applying it at
\[
\theta=\Theta(\omega),\qquad \theta'=\Theta'(\omega),\qquad
\mu=\Law(\Theta),\qquad \mu'=\Law(\Theta')
\]
and taking expectations gives Assumption~\ref{ass:mono}, so the
pointwise condition is the stronger one. The converse
fails: the expectation may average sign-indefinite pointwise pairings into a dissipative one
(Proposition~\ref{prop:strict}). Example~\ref{ex:concrete}(i) verifies
Assumption~\ref{ass:mono} exactly, Example~\ref{ex:fulllaw} perturbatively.
\end{remark}

\begin{remark}[Coupling formulation]\label{rem:intrinsic}
For $t>0$ the $\sigma$-field $\F_t$ supports a uniformly distributed variable (a
function of $W^1_t$), so every coupling of two laws in $\mathcal P_2(\mathcal H)$ is
the joint law of a pair of $\F_t$-measurable
inputs, by the kernel representation lemma~\cite[Lemma~4.22]{Kallenberg}. The diagonal hypotheses,
Assumptions~\ref{ass:lip} and~\ref{ass:mono}, are therefore intrinsic: properties of
the coefficients, not of the standing basis. As the interior conditions are imposed
$\diff t$-a.e.\ and the terminal ones at $T>0$, no richness assumption on $\F_0$ is
needed. The proofs use only the expected diagonal form.
\end{remark}
\begin{remark}[The $\beta_1$ dichotomy in the jump channel]\label{rem:beta1-dich}
Taking $\delta X=\delta Y=\delta Z=0$ and $\delta U$ arbitrary in \eqref{eq:mono} reduces the
interior pairing to
\[
\E\int_E\inner{\delta h(e)}{\Gstar\delta U(e)}\,\nu(\diff e)
\le-\beta_1\,\E\Lnorm{\Gstar\delta U}^2 ,
\]
so $\beta_1>0$ forces the jump coefficient $h$ to be strictly $U$-dissipative. A jump
coefficient of the natural additive form $h=h(t,x,e)$ carries no $U$-feedback, so
$\delta h=0$ while $\delta U$ may be chosen with $\Gstar\delta U\neq0$; the displayed bound
then forces $\beta_1\le0$, hence $\beta_1=0$ as the constants are non-negative.
Assumption~\ref{ass:mono} then requires $m\ge n$ with $\alpha,\beta_2>0$, the
$G$-injective branch, and the $m<n$ rank case is available only for $U$-dissipative jump
coefficients.
The restriction is structural: it follows from Assumption~\ref{ass:mono} itself, not from
the estimates that exploit it. Recovering the additive form at $m<n$ under a different
structural condition remains open.
\end{remark}
\begin{remark}[Dissipativity margin]\label{rem:margin}
Say that \eqref{eq:mono} holds \emph{with margin} $c>0$ if
\begin{equation}\label{eq:monomargin}
\begin{aligned}
&\E\!\Big[\inner{\delta b}{\Gstar\delta Y}+\inner{\delta\sigma}{\Gstar\delta Z}_F
+\inner{-\delta f}{G\delta X}+\!\int_E\!\inner{\delta h(e)}{\Gstar\delta U(e)}\nu(\diff e)\Big]\\
&\quad\le-\beta_1\E\big[|\Gstar\delta Y|^2+|\Gstar\delta Z|_F^2+\Lnorm{\Gstar\delta U}^2\big]-\beta_2\E|G\delta X|^2\\
&\qquad\ \ -c\,\E\big[|\Gstar\delta Y|^2+|\Gstar\delta Z|_F^2+\Lnorm{\Gstar\delta U}^2+|G\delta X|^2\big] .
\end{aligned}
\end{equation}
Suppose the interior pairing has been bounded with constants $\bar\beta_1,\bar\beta_2$. For
every $0<c<\min\{\bar\beta_1,\bar\beta_2\}$ the same bound is \eqref{eq:monomargin} with the
reduced constants $\beta_i:=\bar\beta_i-c$ and margin $c$: recording weaker constants turns the
difference into slack, which a perturbation raising the interior pairing by at most the margin
leaves intact. When $m=n$ the margin may instead be recorded in the \emph{unweighted} form
\begin{equation}\label{eq:monomarginu}
\begin{aligned}
&\E\!\Big[\inner{\delta b}{\Gstar\delta Y}+\inner{\delta\sigma}{\Gstar\delta Z}_F
+\inner{-\delta f}{G\delta X}+\!\int_E\!\inner{\delta h(e)}{\Gstar\delta U(e)}\nu(\diff e)\Big]\\
&\quad\le-\beta_1\E\big[|\Gstar\delta Y|^2+|\Gstar\delta Z|_F^2+\Lnorm{\Gstar\delta U}^2\big]-\beta_2\E|G\delta X|^2\\
&\qquad\ \ -c\,\E\big[|\delta Y|^2+|\delta Z|_F^2+\Lnorm{\delta U}^2+|\delta X|^2\big] .
\end{aligned}
\end{equation}
\end{remark}
